\begin{proposition}[负实根主导的二周期比值误差：共振区 \texorpdfstring{$R_m$}{Rm} 的交错逼近（可复核）]\label{prop:pom-resonance-two-cycle-ratio}
\leanverified{paper\_pom\_resonance\_two\_cycle\_ratio}
设 \(\{S(m)\}_{m\ge 0}\) 满足某个 \(d\) 阶整系数线性递推，其特征多项式 \(P(\lambda)\) 的根按模排序为
\[
|\lambda_1|>|\lambda_2|\ge|\lambda_3|\ge\cdots\ge|\lambda_d|.
\]
并假设：
\begin{enumerate}
  \item \(\lambda_1>0\) 为单根；
  \item \(\lambda_2<0\) 为实根；
  \item 解的谱分解中，\(\lambda_1\) 与 \(\lambda_2\) 的模态系数均非零（一般位置；对本文的 \(S_q\) 数值上成立）。
\end{enumerate}
令比值序列
\[
R_m:=\frac{S(m+1)}{S(m)}.
\]
则存在常数 \(A\neq 0\)，使得当 \(m\to\infty\) 时有渐近展开
\[
\boxed{\;
R_m=\lambda_1\ +\ A\left(\frac{\lambda_2}{\lambda_1}\right)^{m}\ +\ O\!\left(\left(\max\!\left\{\frac{|\lambda_3|}{\lambda_1},\left(\frac{|\lambda_2|}{\lambda_1}\right)^2\right\}\right)^{m}\right)\;}
\]
从而主误差项携带一个\textbf{强二周期交替}：\(R_m-\lambda_1\) 在大 \(m\) 上呈现符号振荡（不再单调收敛）。
\end{proposition}
\begin{proof}
线性递推解的一般理论给出（在 \(\lambda_1\) 为单根且严格主导时）存在常数 \(C_1,C_2\neq 0\) 使
\[
S(m)=C_1\lambda_1^m+C_2\lambda_2^m+O(|\lambda_3|^m).
\]
令 \(\rho:=\lambda_2/\lambda_1\) 与 \(\delta:=|\lambda_3|/\lambda_1\)，并把上式代入
\(
R_m=S(m+1)/S(m)
\)：
\[
\frac{R_m}{\lambda_1}
=\frac{1+\frac{C_2}{C_1}\rho^{m+1}+O(\delta^{m+1})}{1+\frac{C_2}{C_1}\rho^{m}+O(\delta^{m})}
=1+\frac{C_2}{C_1}(\rho-1)\rho^{m}+O(\delta^m)+O(\rho^{2m}).
\]
取 \(A:=\lambda_1\frac{C_2}{C_1}(\rho-1)\neq 0\) 即得结论，并且由于 \(\rho<0\) 从而 \(\rho^m\) 交替变号，主误差为二周期交替项。
\end{proof}

\begin{corollary}[Aitken \texorpdfstring{$\Delta^2$}{Delta^2} 的严格作用：消去二周期主几何项]\label{cor:pom-aitken-cancels-two-cycle}
\leanverified{paper\_pom\_aitken\_cancels\_two\_cycle}
对任意数列 \(\{x_m\}\) 定义 Aitken 变换
\[
\mathcal{A}(x_m):=x_m-\frac{(x_{m+1}-x_m)^2}{x_{m+2}-2x_{m+1}+x_m}.
\]
在命题 \ref{prop:pom-resonance-two-cycle-ratio} 的假设下，
\[
\boxed{\;
\mathcal{A}(R_m)=\lambda_1\ +\ O\!\left(\left(\max\!\left\{\frac{|\lambda_3|}{\lambda_1},\left(\frac{|\lambda_2|}{\lambda_1}\right)^2\right\}\right)^{m}\right)\;}
\]
因此当负实根 \(\lambda_2\) 为次主导模态时，Aitken \(\Delta^2\) 并非经验技巧，而是一个\textbf{严格的二周期误差消去器}：它把由 \(\lambda_2\) 诱导的交替主误差项从比值估计中消去，把收敛阶提升到由下一层谱尺度控制。
\end{corollary}
\begin{proof}
把命题 \ref{prop:pom-resonance-two-cycle-ratio} 写为
\(
R_m=\lambda_1+a\rho^{m}+O(\tau^{m})
\)
其中 \(a\neq 0\)、\(\rho=\lambda_2/\lambda_1\) 且
\(
\tau:=\max\{\delta,|\rho|^2\}<|\rho|<1
\)。
直接代入差分算子 \(\Delta\) 与 \(\Delta^2\) 得
\[
\Delta R_m=a(\rho-1)\rho^{m}+O(\tau^{m}),
\qquad
\Delta^2 R_m=a(\rho-1)^2\rho^{m}+O(\tau^{m}).
\]
于是
\[
\frac{(\Delta R_m)^2}{\Delta^2 R_m}
=a\rho^{m}+O(\tau^{m}),
\]
从而
\(
\mathcal{A}(R_m)=R_m-a\rho^{m}+O(\tau^{m})
=\lambda_1+O(\tau^{m})
\)。
\end{proof}

\begin{definition}[共振谱—算法指纹：\texorpdfstring{$\sigma_{\mathrm{osc}}(q)$}{sigma\_osc(q)} 与 \texorpdfstring{$\delta_{\mathrm{ait}}(q)$}{delta\_ait(q)}]\label{def:pom-resonance-aitken-fingerprints}
对共振区的特征多项式 \(P_q(\lambda)\) 的根按模排序
\[
|\lambda_{1,q}|>|\lambda_{2,q}|\ge|\lambda_{3,q}|\ge\cdots,
\qquad \lambda_{1,q}=r_q,
\]
定义两条\textbf{可直接由多项式求根得到}的共振指纹：
\[
\boxed{\;
\sigma_{\mathrm{osc}}(q):=\frac{|\lambda_{2,q}|}{r_q},
\qquad
\delta_{\mathrm{ait}}(q):=\frac{|\lambda_{3,q}|}{r_q}\;}
\]
其中 \(\sigma_{\mathrm{osc}}\) 控制比值 \(R_m\) 的二周期主误差强度，而 \(\delta_{\mathrm{ait}}\) 控制 Aitken \(\Delta^2\) 后的主误差强度（推论 \ref{cor:pom-aitken-cancels-two-cycle}）。
\end{definition}

\begin{remark}[窗口长度的审计证书：用 \texorpdfstring{$(\sigma_{\mathrm{osc}},\delta_{\mathrm{ait}})$}{(sigma\_osc,delta\_ait)} 预测有限窗收敛]\label{rem:pom-aitken-window-audit}
在一般位置下，存在常数 \(C,C'>0\) 使得对充分大的 \(m\),
\[
\left|R_m-r_q\right|\le C\,\sigma_{\mathrm{osc}}(q)^{m},
\qquad
\left|\mathcal{A}(R_m)-r_q\right|\le C'\,\bigl(\max\{\delta_{\mathrm{ait}}(q),\sigma_{\mathrm{osc}}(q)^2\}\bigr)^{m}.
\]
因此若目标相对误差阈值为 \(\mathrm{tol}\in(0,1)\)，则一个直接的窗口长度准则是选取
\[
m\ \gtrsim\ \frac{\log(\mathrm{tol})}{\log \delta_{\mathrm{ait}}(q)}
\qquad(\text{Aitken 加速后；若不用 Aitken 则用 }\sigma_{\mathrm{osc}}(q)\text{ 代替}).
\]
对本文 \(q=9,\dots,17\) 的共振多项式，按模排序的次主导根在数值上稳定为负实根，从而满足二周期交替结构；相应的 \(\sigma_{\mathrm{osc}}(q)\) 与 \(\delta_{\mathrm{ait}}(q)\) 由脚本 \texttt{scripts/exp\_fold\_collision\_resonance\_aitken\_certificate\_q9\_17.py} 一键生成，见表 \ref{tab:fold_collision_resonance_aitken_certificate_q9_17}。
并且在该范围内数值上有 \(\delta_{\mathrm{ait}}(q)>\sigma_{\mathrm{osc}}(q)^2\)，故推论 \ref{cor:pom-aitken-cancels-two-cycle} 中的 \(\max\{\delta_{\mathrm{ait}},\sigma_{\mathrm{osc}}^2\}\) 在该窗内可简化为 \(\delta_{\mathrm{ait}}\)。
在该范围内两者随 \(q\) 增大呈现上升趋势，意味着在固定 \(m\)-窗口下收敛会系统性变慢，而 Aitken 的收益会随之更加显著。
\end{remark}

\begin{corollary}[算法性推论：Aitken 窗长的“黄金不变性”]\label{cor:pom-aitken-window-golden-invariance}
\leanverified{paper\_pom\_aitken\_window\_golden\_invariance}
沿用注 \ref{rem:pom-aitken-window-audit} 的窗口长度准则。
若在共振饱和极限中 Aitken 主误差强度满足
\[
\delta_{\mathrm{ait}}(q)\ \longrightarrow\ \varphi^{-1}
\qquad(q\to\infty)
\]
（例如推论 \ref{cor:pom-aitken-golden-limit} 的条件成立），
则对任意固定相对误差阈值 \(\mathrm{tol}\in(0,1)\)，Aitken 加速所需的极限窗长与 \(q\) 脱钩，并满足
\[
m\ \sim\ \frac{\log(1/\mathrm{tol})}{\log\varphi}
\qquad(q\to\infty).
\]
因此在“两步饱和 \(\eta_q\to 1\)” 的共振极限中，Aitken 并非经验加速，而是把收敛常数刚性化为黄金比：窗口复杂度退化为一个由 \(\varphi\) 唯一控制的对数尺度。
\end{corollary}
\begin{proof}
由注 \ref{rem:pom-aitken-window-audit} 的准则，
对充分大 \(m\) 取
\(
m\gtrsim \frac{\log(\mathrm{tol})}{\log\delta_{\mathrm{ait}}(q)}
\)。
在 \(\delta_{\mathrm{ait}}(q)\to\varphi^{-1}\) 下，
\(
\log\delta_{\mathrm{ait}}(q)\to -\log\varphi
\)，
从而
\(
\frac{\log(\mathrm{tol})}{\log\delta_{\mathrm{ait}}(q)}\to \frac{\log(1/\mathrm{tol})}{\log\varphi}
\)；
即得极限窗长的渐近式。证毕。
\end{proof}

\begin{table}[H]
\centering
\scriptsize
\setlength{\tabcolsep}{4pt}
\caption{Aitken-certificate root fingerprints for resonance-region Fold collision moment characteristic polynomials $P_q$ (audited recurrences, $q=9,\dots,17$). Roots are sorted by modulus: $\lambda_1=r_q$ (Perron root), $\lambda_2$ (2nd), $\lambda_3$ (3rd). We report $\sigma_{\mathrm{osc}}:=|\lambda_2|/r_q$ (dominant alternating ratio-error strength) and $\delta_{\mathrm{ait}}:=|\lambda_3|/r_q$ (post-Aitken error strength).}
\label{tab:fold_collision_resonance_aitken_certificate_q9_17}
\begin{tabular}{r r r r r r r}
\toprule
$q$ & order $d_q$ & $r_q$ & $\lambda_2$ & $\sigma_{\mathrm{osc}}$ & $|\lambda_3|$ & $\delta_{\mathrm{ait}}$\\
\midrule
9 & 7 & 11.778422 & -7.065331 & 0.5999 & 5.807034 & 0.4930\\
10 & 9 & 14.771098 & -9.100142 & 0.6161 & 7.686459 & 0.5204\\
11 & 9 & 18.535847 & -11.712939 & 0.6319 & 10.074255 & 0.5435\\
12 & 13 & 23.273460 & -15.066709 & 0.6474 & 13.109656 & 0.5633\\
13 & 11 & 29.237339 & -19.369971 & 0.6625 & 16.968007 & 0.5804\\
14 & 13 & 36.747376 & -24.889362 & 0.6773 & 21.870525 & 0.5952\\
15 & 11 & 46.207510 & -31.965663 & 0.6918 & 28.096490 & 0.6081\\
16 & 13 & 58.127952 & -41.034237 & 0.7059 & 35.998554 & 0.6193\\
17 & 13 & 73.153329 & -52.651060 & 0.7197 & 46.021996 & 0.6291\\
\bottomrule
\end{tabular}
\end{table}

